\documentclass[12pt, a4paper]{article}
\usepackage[utf8x]{inputenc}
\usepackage[T2A]{fontenc}
\usepackage[russian]{babel}
\usepackage{amsmath}
\usepackage{amssymb}
\usepackage{graphicx}
\usepackage{hyperref}
\usepackage{enumitem}

\setlist{nosep, leftmargin=*}

\usepackage{titling}

\usepackage{geometry}

\geometry{a4paper, top=2cm, bottom=2cm, left=2.5cm, right=1.5cm}

\begin{document}

\begin{titlepage}
    \begin{center}
        \vspace*{6cm}

        \textbf{\Huge Реферат}\\[0.5cm]
        \textbf{\Large на тему: «Тензор инерции (вывод).}\\[0.2cm]
        \textbf{\Large Эффект Джанибекова»}\\[1cm]
        \textbf{\Large Вариант 13}\\[3cm]        
        \textbf{\Large Поповкин Артемий Андреевич}\\[0.3cm]
        \textbf{\large Б9122-02.03.01сцт}\\[9cm]

        \textsc{\Large 2026}

        \vfill
    \end{center}
\end{titlepage}

\section{Тензор инерции: вывод и свойства}

\subsection{Момент импульса твердого тела}
Рассмотрим твердое тело как систему материальных точек с массами $m_i$ и радиус-векторами $\mathbf{r}_i$ в неподвижной системе координат. Введем систему отсчета, жестко связанную с телом, с началом в точке $O$ (полюсе). Абсолютная скорость точки:
\[
\mathbf{v}_i = \mathbf{v}_O + \vec{\omega} \times \mathbf{r}_i',
\]
где $\vec{\omega}$ — мгновенная угловая скорость вращения твёрдого тела, $\mathbf{r}_i'$ — радиус-вектор точки относительно полюса $O$ в связанной системе.

Момент импульса относительно полюса $O$:
\[
\mathbf{L}_O = \sum_i \mathbf{r}_i' \times m_i \mathbf{v}_i.
\]
Подставляя $\mathbf{v}_i$ и полагая $\mathbf{v}_O = 0$ (полюс покоится или совпадает с центром масс), получаем:
\[
\mathbf{L}_O = \sum_i m_i \mathbf{r}_i' \times (\vec{\omega} \times \mathbf{r}_i') = \sum_i m_i \left[ \vec{\omega} (\mathbf{r}_i' \cdot \mathbf{r}_i') - \mathbf{r}_i' (\mathbf{r}_i' \cdot \boldsymbol{\omega}) \right].
\]

\subsection{Определение тензора инерции}
В компонентной записи (индексы $k, l, m = 1,2,3$):
\[
L_k = \sum_i m_i \left( \delta_{km} r_i^2 - x_{ik} x_{im} \right) \omega_m = I_{km} \omega_m,
\]
где $r_i^2 = x_{i1}^2 + x_{i2}^2 + x_{i3}^2$, $\delta_{km}$ - символ Кронекера $ = \begin{cases}
    1, & k = m \\
    0, & k  \neq m
\end{cases}$.

Тензор инерции второго ранга:
\[
I_{km} = \sum_i m_i \left( r_i^2 \delta_{km} - x_{ik} x_{im} \right)
\]
Для непрерывного распределения массы с плотностью $\rho(\mathbf{r})$:
\[
I_{km} = \int_V \rho(\mathbf{r}) \left( r^2 \delta_{km} - x_k x_m \right) dV
\]

\subsection{Компоненты тензора}
В декартовых координатах \((x, y, z)\) компоненты тензора инерции имеют вид:
\begin{itemize}
\item \textbf{Диагональные} (моменты инерции относительно осей):
  \[
  I_{xx} = \int (y^2 + z^2)\,dm,\quad
  I_{yy} = \int (x^2 + z^2)\,dm,\quad
  I_{zz} = \int (x^2 + y^2)\,dm.
  \]
\item \textbf{Недиагональные} (центробежные моменты инерции):
  \[
  I_{xy} = -\int xy\,dm,\quad
  I_{xz} = -\int xz\,dm,\quad
  I_{yz} = -\int yz\,dm.
  \]
\end{itemize}

\textbf{Пример 1.} Две точечные массы \(m_1 = m\) и \(m_2 = m\) расположены в точках \((a,0,0)\) и \((0,a,0)\) соответственно. Вычислим компоненты тензора инерции относительно начала координат.

$$
\begin{array}{ll}
I_{xx} = \sum m_i (y_i^2 + z_i^2) = m(0^2+0^2) + m(a^2+0^2) = m a^2 \\
I_{yy} = \sum m_i (x_i^2 + z_i^2) = m(a^2+0^2) + m(0^2+0^2) = m a^2 \\
I_{zz} = \sum m_i (x_i^2 + y_i^2) = m(a^2+0^2) + m(0^2 + a^2) = 2 m a^2 \\
I_{xy} = -\sum m_i x_i y_i = -(m \cdot a \cdot 0 + m \cdot 0 \cdot a) = 0
\end{array}
$$


Аналогично \(I_{xz}=0\), \(I_{yz}=0\).
Тензор диагоналей: \(\mathbf{I} = \begin{pmatrix} m a^2 & 0 & 0 \\ 0 & m a^2 & 0 \\ 0 & 0 & 2 m a^2 \end{pmatrix}\).

\textbf{Пример 2.} Тонкий стержень длины \(L\) и массы \(M\), расположенный вдоль оси \(x\) от \(-L/2\) до \(+L/2\). Плотность \(\lambda = M/L\). Вычислим \(I_{yy}\) (момент инерции относительно оси \(y\)):
\[
I_{yy} = \int_{-L/2}^{L/2} (x^2 + z^2)\, dm = \int_{-L/2}^{L/2} x^2 \cdot \lambda\, dx = \lambda \left[ \frac{x^3}{3} \right]_{-L/2}^{L/2} = \frac{M}{L} \cdot \frac{L^3}{12} = \frac{1}{12} M L^2.
\]
Недиагональные компоненты (например, \(I_{xy}\)) для стержня на оси \(x\) равны нулю, так как \(y=0\).

\subsection{Свойства тензора}
\begin{enumerate}
    \item \textbf{Симметричность}: $I_{km} = I_{mk}$.
    \item \textbf{Положительная определенность}: для любого ненулевого $\boldsymbol{\omega}$ кинетическая энергия вращения 
    \[
    T = \frac{1}{2} I_{km} \omega_k \omega_m > 0
    \]
    \item \textbf{Преобразование при повороте координат}: компоненты преобразуются как тензор второго ранга:
    \[
    I'_{pq} = a_{pi} a_{qj} I_{ij}
    \]
    где $a_{pi}$ — матрица поворота.
    \item \textbf{Существование главных осей}: вследствие симметричности тензора существует ортонормированный базис, в котором тензор имеет диагональный вид:
    \[
    \mathbf{I} = \begin{pmatrix}
    I_1 & 0 & 0 \\
    0 & I_2 & 0 \\
    0 & 0 & I_3
    \end{pmatrix}
    \]
    где $I_1, I_2, I_3$ — главные моменты инерции. Оси тензора (ортонормированный базис) называются главными осями инерции. Для тела с произвольной формой главные оси можно найти, решив задачу на собственные значения $I \mathbf{e} = \lambda \mathbf{e}$.
\end{enumerate}

\section{Эффект Джанибекова (теорема о теннисной ракетке)}

\subsection{Постановка задачи}
Эффект Джанибекова — это явление неустойчивости свободного вращения твердого тела вокруг оси, соответствующей среднему по величине главному моменту инерции. Впервые оно было замечено космонавтом В. Джанибековым в невесомости, но математически описано как часть решения уравнений Эйлера для твердого тела.

Рассмотрим твердое тело, закрепленное в центре масс (или свободное в невесомости), на которое не действуют внешние моменты сил. Уравнения движения в главных осях (уравнения Эйлера) имеют вид:
\[
\begin{cases}
I_1 \dot{\omega}_1 = (I_2 - I_3) \omega_2 \omega_3, \\
I_2 \dot{\omega}_2 = (I_3 - I_1) \omega_3 \omega_1, \\
I_3 \dot{\omega}_3 = (I_1 - I_2) \omega_1 \omega_2.
\end{cases}
\]

\subsection{Анализ устойчивости стационарных вращений}
Стационарное вращение соответствует $\dot{\boldsymbol{\omega}} = 0$. Пусть, например, $\omega_1 = \Omega, \omega_2 = \omega_3 = 0$ (вращение вокруг главной оси 1). Исследуем устойчивость, задав малые возмущения: $\omega_2, \omega_3 \ll \Omega$. Линеаризуем уравнения:
\[
I_2 \dot{\omega}_2 = (I_3 - I_1) \Omega \omega_3, \quad I_3 \dot{\omega}_3 = (I_1 - I_2) \Omega \omega_2
\]
Продифференцируем первое и подставим второе:
\[
\ddot{\omega}_2 = \frac{(I_3 - I_1)(I_1 - I_2)}{I_2 I_3} \Omega^2 \omega_2
\]
Обозначим $\lambda = \dfrac{(I_3 - I_1)(I_1 - I_2)}{I_2 I_3} \Omega^2$.

Устойчивость требует $\lambda < 0$ (колебательный режим).
Рассмотрим три случая упорядочения моментов инерции: $I_1 < I_2 < I_3$.

\begin{enumerate}
    \item \textbf{Вращение вокруг оси 1 (наименьший момент)}.  
    $I_1 - I_2 < 0,\ I_3 - I_1 > 0 \Rightarrow$ произведение $(I_3 - I_1)(I_1 - I_2) < 0 \Rightarrow \lambda < 0 \Rightarrow$ \textbf{устойчиво}.
    
    \item \textbf{Вращение вокруг оси 3 (наибольший момент)}.  
    $I_3 - I_1 > 0,\ I_1 - I_2 < 0 \Rightarrow$ произведение $(I_3 - I_1)(I_1 - I_2) < 0 \Rightarrow \lambda < 0 \Rightarrow$ \textbf{устойчиво}.
    
    \item \textbf{Вращение вокруг оси 2 (средний момент)}.  
    $I_2 - I_1 > 0,\ I_3 - I_2 > 0 \Rightarrow$ произведение $(I_3 - I_2)(I_2 - I_1) > 0 \Rightarrow \lambda > 0 \Rightarrow$ \textbf{неустойчиво}.
\end{enumerate}

Таким образом, вращение вокруг оси со средним главным моментом инерции является неустойчивым: малые возмущения экспоненциально нарастают, приводя к периодическому «кувырканию» тела.

\subsection{Геометрическая интерпретация}
Эффект можно наглядно представить с помощью эллипсоида инерции. В отсутствие моментов сохраняются:
\begin{itemize}
    \item Кинетическая энергия $T = \dfrac{1}{2} (I_1 \omega_1^2 + I_2 \omega_2^2 + I_3 \omega_3^2)$;
    \item Квадрат момента импульса $L^2 = I_1^2 \omega_1^2 + I_2^2 \omega_2^2 + I_3^2 \omega_3^2$.
\end{itemize}
Траектория конца вектора $\boldsymbol{\omega}$ лежит на пересечении эллипсоида энергии (эллипсоид Пуансо) и сферы момента. При $I_1 < I_2 < I_3$ точки пересечения, соответствующие вращению вокруг оси 2, лежат на сепаратрисах, разделяющих области устойчивых движений вокруг осей 1 и 3. Любое сколь угодно малое отклонение приводит к переходу от одной устойчивой оси к другой.

\subsection{Демонстрация и практическое значение}
В условиях невесомости (МКС) эффект проявляется как периодические перевороты вращающегося предмета (гайка, ракетка) на $180^\circ$. Период переворота зависит от разности моментов инерции и начальной угловой скорости. Этот эффект необходимо учитывать при проектировании космических аппаратов и ориентации спутников, чтобы избежать неконтролируемого вращения.

\end{document}